\chapter{The Ontology of Reflexive Constraint Realism}
\label{ch:ontology}

\begin{chapterguide}
This chapter states the ontological side of Reflexive Constraint Realism
(\RCR). It defines reality, entities, processes, structures, laws, levels, and
social facts, then gives the framework's axioms and graded commitment ladder.
The aim is a defensible realism that does not pretend that science supplies a
complete inventory of being.
\end{chapterguide}

\RCR\ begins with a modest but non-negotiable ontological claim: the world does
not depend on human inquiry for its existence. Stars existed before astronomy.
Biological reproduction preceded genetics. People and institutions can affect
one another before social scientists classify them. A mind-independent world is
not necessarily a world described by one final vocabulary. It is a world whose
behaviour can resist, frustrate, and correct our descriptions.

The word \emph{constraint}\index{constraint} is central. A constraint is a
stable restriction on what can happen, or a stable tendency that makes some
changes more likely than others, under specified conditions. A wall constrains
motion. A chemical bond constrains molecular configurations. A legal rule
constrains behaviour through sanctions and expectations. A statistical
relationship can express a constraint when it remains stable under relevant
changes. The concept is broad enough to travel across sciences but precise
enough to require a scope.

\section{Generative reality}

The ontology of \RCR\ is \emph{generative}: reality contains processes that
produce events, not only a collection of regular sequences. A mechanism, field,
institution, or ecological relation matters because it makes a difference to
what can happen. Generativity does not require determinism. A stochastic
process can generate a distribution. A social institution can make behaviour
more likely without forcing each individual action.

The ontology is also \emph{layered}. Physical processes, organisms, persons,
institutions, and ecosystems can be dependent on one another without being
interchangeable descriptions. A higher-level entity is real when it has
stable powers, boundaries, or relations that make a difference to outcomes,
even if it is physically realized by lower-level processes.

Finally, the ontology is \emph{open}. Science may discover new kinds of entity,
and current categories may be replaced. The framework should therefore state
what warrants commitment without claiming to know the final catalogue of
reality.

\section{Eight axioms}

The following axioms are methodological-metaphysical starting points. They are
not self-evident truths from which all science can be deduced. They are
commitments that make the project of objective inquiry intelligible and expose
where it can fail.

\begin{principlebox}
\textbf{Axiom 1: Mind-independent generativity.} There is a reality whose
processes do not depend for their existence on being represented by human
agents. Those processes generate resistances, regularities, possibilities, and
events that inquiry can discover imperfectly.
\end{principlebox}

\begin{principlebox}
\textbf{Axiom 2: Constraint plurality.} The world constrains inquiry through
more than one kind of relation: material interaction, causal production,
mathematical structure, temporal dependence, ecological organization, and
social rule. No single layer is assumed to exhaust all scientifically relevant
being.
\end{principlebox}

\begin{principlebox}
\textbf{Axiom 3: Mediated access.} Scientific access to reality is always
mediated by concepts, instruments, models, language, and institutions. Mediation
does not imply invention; it creates a route through which independent
constraints can act on claims.
\end{principlebox}

\begin{principlebox}
\textbf{Axiom 4: Fallibility.} Every finite scientific claim is revisable.
Confidence can be high without being certain, and ontological commitment can
be strong without being final.
\end{principlebox}

\begin{principlebox}
\textbf{Axiom 5: Invariance matters.} A claim gains epistemic and ontological
support when its relevant consequences remain stable under admissible changes
of observer, instrument, representation, intervention, or context.
\end{principlebox}

\begin{principlebox}
\textbf{Axiom 6: Situated inquiry.} Every inquiry begins from a location in a
social, material, and historical system. The location affects what becomes
visible, but it need not determine what is true.
\end{principlebox}

\begin{principlebox}
\textbf{Axiom 7: Public corrigibility.} Objective inquiry requires practices
through which competent others can inspect, challenge, reproduce, and revise a
claim. Private certainty is not a substitute for public correction.
\end{principlebox}

\begin{principlebox}
\textbf{Axiom 8: Scope-bound truth.} A claim is true only with respect to the
objects, variables, scales, interventions, and background conditions that its
evidence supports. Generality must be earned by transport.
\end{principlebox}

Axioms 1 and 3 must be held together. If one accepts only Axiom 1, one risks
naive realism: the world is real, so observations are treated as transparent.
If one accepts only Axiom 3, one risks constructivist relativism: all access is
mediated, so reality becomes indistinguishable from representation. Axioms 5
through 8 explain how mediation can be disciplined.

\section{What scientific entities are}

An entity is not simply whatever appears as a noun in a theory. \RCR\ uses
three tests.

\begin{enumerate}
\item \textbf{Boundary test:} Can the entity be distinguished from neighbouring
processes by a stable criterion?
\item \textbf{Difference test:} Does changing or removing the entity make a
reliable difference to outcomes?
\item \textbf{Convergence test:} Do independent measurements, models, or
interventions support the entity's role?
\end{enumerate}

An entity can pass these tests without being fundamental. A gene may be a
molecular sequence, a regulatory function, or a population-level unit depending
on the question. A firm is not a biological organism, but it can be a real
social entity if it has legally recognized boundaries, resources, decision
procedures, and causal effects.

\RCR\ therefore distinguishes \emph{realization} from \emph{reduction}. To say
that a higher-level entity is realized by lower-level processes is to say that
the higher-level entity depends on them. To say that it is reducible is to say
that its explanatory and causal role can be replaced without loss. The first
does not entail the second.

\section{Structures and relations}

A structure consists of elements together with relations and transformations.
Scientific theories often support relations more strongly than they support
the intrinsic nature of elements. A model may establish that changes in \(X\)
alter \(Y\) with a certain pattern without determining whether \(X\) is a
fundamental property or a useful aggregate.

The realist commitment to structure must be disciplined by relevance. A
mathematical isomorphism alone does not show that two physical systems share
the same ontology. The relation must survive a mapping between model and
target, and it must matter to prediction, intervention, or explanation.

\section{Laws as modal constraints}

\RCR\ treats laws of nature as candidates for stable modal constraints: they
describe what would happen under a range of admissible changes, not merely what
has happened so far. A law-like claim can be probabilistic, approximate, or
scale-specific. It becomes law-like when its counterfactual and interventional
consequences are stable enough to guide inquiry.

This view avoids the choice between universal regularity and mere description.
A law can be a real feature of a domain without being the only law that exists
or without applying outside the regime in which its variables are well-defined.

\section{The ontology of social facts}

Social facts are not less real because they depend on human activity. Money,
property, citizenship, corporations, and scientific institutions have histories
and rules, but they constrain behaviour and distribute resources. Their
existence is relational and reflexive: they persist because people recognize,
enforce, or enact them, yet once established they are not freely adjustable by
each individual.

The ontology of social science must therefore include at least:

\begin{itemize}
\item material bodies and environments;
\item agents with capacities and interpretations;
\item rules, norms, and institutions;
\item networks of dependence and power;
\item historical trajectories that alter future possibilities.
\end{itemize}

A social explanation should not collapse a norm into a molecule, but neither
should it detach the norm from the people, artifacts, and sanctions that
realize it.

\section{Graded ontological commitment}

\RCR\ recommends a commitment ladder. The ladder is not a probability scale.
It records what sort of support has been earned.

\begin{table}[ht]
\centering
\smalltable
\begin{tabularx}{\textwidth}{p{0.25\textwidth}Y Y}
\toprule
Level & Commitment & Typical warrant\\
\midrule
L0: useful posit & Use the term without belief in its independent reality & Predictive or computational utility\\
L1: empirical pattern & Believe that a stable phenomenon occurs & Convergent measurement and replication\\
L2: structural commitment & Believe that a relation or invariance is real in a domain & Cross-model stability and limiting relations\\
L3: causal commitment & Believe that an entity, process, or mechanism makes a difference & Intervention, counterfactual robustness, mechanism\\
L4: domain ontology & Treat a family of entities and relations as real in a mature domain & Long-run convergence, transport, integration, and correction\\
\bottomrule
\end{tabularx}
\caption{The \RCR\ commitment ladder. Movement upward requires stronger and
more independent constraints; movement downward remains possible after
counterevidence.}
\label{tab:commitment-ladder}
\end{table}

The ladder protects against two errors. It prevents a useful model from being
treated as a literal inventory. It also prevents anti-realism from ignoring
the epistemic significance of interventions that could not succeed if the
relevant entity or structure were absent.

\section{A layered map of reality}

\begin{figure}[ht]
\centering
\begin{tikzpicture}[
  layer/.style={draw=linegray,rounded corners=2pt,align=center,minimum width=5.0cm,minimum height=0.72cm,font=\small},
  ar/.style={-{Latex[length=2mm]},draw=teal,thick}
]
\node[layer,fill=covernavy!12] (phys) {Physical processes and fields};
\node[layer,fill=teal!12,above=4mm of phys] (bio) {Organisms, regulation, ecology};
\node[layer,fill=teal!8,above=4mm of bio] (agent) {Agents, cognition, interpretation};
\node[layer,fill=warm,above=4mm of agent] (soc) {Institutions, norms, power, history};
\draw[ar] (phys) -- (bio);
\draw[ar] (bio) -- (agent);
\draw[ar] (agent) -- (soc);
\draw[ar,dashed] (soc.east) to[out=0,in=0] node[right,font=\scriptsize,align=left]{feedback\\and selection} (phys.east);
\end{tikzpicture}
\caption{A non-reductive layered ontology. Higher layers are realized by lower
processes while feeding back through organization, action, and selection.}
\label{fig:ontology-layers}
\end{figure}

The diagram is not a ladder of value or a claim that reality has exactly four
levels. It is a reminder that scientific explanations can move across levels.
A social institution is materially realized, psychologically interpreted, and
historically reproduced. A disease is biological, embodied, measured, and
institutionally classified. Objectivity depends on preserving the causal links
between levels rather than pretending that one level can answer every question.

\section*{Chapter summary}

\RCR\ posits a mind-independent, generative, layered, and open reality. Its
basic unit is the constraint: a stable restriction or tendency that survives
specified changes. Entities earn commitment through boundary, difference, and
convergence tests. Structures are real when relevant relations remain stable
across representations and interventions. Laws are modal, scope-bound
constraints. Social facts are real relational entities. The eight axioms and
the commitment ladder make realism graded, fallible, and accountable.

\begin{discussion}
\begin{enumerate}
\item Can an entity be real without being fundamental?
\item What would count as evidence that a higher-level social mechanism has
causal powers?
\item Does the commitment ladder describe degrees of belief, degrees of truth,
or degrees of warrant?
\end{enumerate}
\end{discussion}

\begin{furtherreading}
Compare the layered approach here with Chakravartty (2007), Cartwright (1999),
Craver (2007), Batterman (2002), Oppenheim and Putnam (1958), and the
literature on reduction and emergence.
\end{furtherreading}

